% !TEX root = task-2-2-1.tex
% Editable task bank. Append new BankTask blocks at the end of the document.
\documentclass[12pt,a4paper]{article}
\usepackage{physics-textbook}
\usepackage{graphicx,needspace}
\usetikzlibrary{decorations.pathmorphing}
\newif\ifShowSolutions
\ifdefined\TutorEdition
  \ShowSolutionstrue
\else
  \ShowSolutionsfalse % Change to \ShowSolutionstrue to display solutions.
\fi
\newcounter{banktask}
\newsavebox{\BankTaskBox}
\newcommand{\BankTaskContent}[3]{%
  \noindent\textbf{\thebanktask.}\quad #2\par
  \ifShowSolutions
    \Needspace{50mm}% Keep the solution heading with its opening text or diagram.
    #3\par
    \ifdefined\TutorEdition
      \ifcsname TutorSource#1\endcsname
        \nopagebreak\smallskip{\small\csname TutorSource#1\endcsname\par}%
      \fi
    \fi
  \fi
}
\newcommand{\BankTask}[3]{%
  \refstepcounter{banktask}\par\medskip
  \sbox{\BankTaskBox}{\begin{minipage}{\linewidth}%
    \BankTaskContent{#1}{#2}{#3}\end{minipage}}%
  \ifdim\dimexpr\ht\BankTaskBox+\dp\BankTaskBox\relax>\dimexpr\textheight-20mm\relax
    \Needspace{35mm}\BankTaskContent{#1}{#2}{#3}%
  \else
    \Needspace{\dimexpr\ht\BankTaskBox+\dp\BankTaskBox\relax}%
    \noindent\usebox{\BankTaskBox}\par
  \fi
  \medskip
}
\begin{document}
\BookSubsection{2.2.1}{Газы}


% Source sdamgia-827: https://phys-ege.sdamgia.ru/problem?id=827
\BankTask{sdamgia-827}{
Какая-либо упорядоченность в расположении частиц вещества отсутствует. Это утверждение соответствует модели

$1$) только газа

$2$) только жидкости

$3$) только твердого тела

$4$) газа, жидкости и твердого тела
}{
\smallskip\noindent\textbf{Решение}\par
Полное отсутствие какой-либо упорядоченности в расположении частиц вещества соответствует только модели газа.

Ответ: $1$.
}

% Source sdamgia-3205: https://phys-ege.sdamgia.ru/problem?id=3205
\BankTask{sdamgia-3205}{
Какое из утверждений не соответствует представлениям молекулярно-кинетической теории о строении газов?

$1$. Все тела состоят из частиц: атомов, молекул и ионов;

$2$. Частицы находятся в непрерывном хаотическом движении (тепловом);

$3$. Все частицы летают со строго определенными по величине скоростями;

$4$. Частицы взаимодействуют друг с другом путем абсолютно упругих столкновений.
}{
\smallskip\noindent\textbf{Решение}\par
Согласно молекулярно-кинетической теории, все тела, состоят из частиц, которые находятся в непрерывном хаотическом движении и взаимодействуют только посредством упругих столкновений. Естественно, что при столкновениях скорости частиц изменяются. Утверждение $3$ не является верным.

Ответ: $3$.
}

% Source sdamgia-3207: https://phys-ege.sdamgia.ru/problem?id=3207
\BankTask{sdamgia-3207}{
Какое из утверждений справедливо для газа?

$1$. Молекулы газа располагаются в строгом порядке

$2$. Газ сохраняет начальный объем

$3$. Характерное расстояние между молекулами газа по порядку величины совпадает с размерами самих молекул

$4$. Газ всегда занимает весь отведенный ему объем
}{
\smallskip\noindent\textbf{Решение}\par
Из всех приведенных утверждений для газа справедливо только последнее. За счет теплового движения, молекулы газа свободно перемещаются по отведенному им объему. Второе утверждение противоречит четвертому. Первое утверждение относится к кристаллическим телам, третье справедливо для твердых тел и жидкостей.

Ответ: $4$.
}

% Source sdamgia-3299: https://phys-ege.sdamgia.ru/problem?id=3299
\BankTask{sdamgia-3299}{
Расстояние между соседними частицами вещества в среднем во много раз превышает размеры самих частиц. Это утверждение соответствует

$1$) только модели строения газов

$2$) только модели строения аморфных тел

$3$) модели строения газов и жидкостей

$4$) модели строения газов, жидкостей и твердых тел
}{
\smallskip\noindent\textbf{Решение}\par
Утверждение о том, что расстояние между соседними частицами вещества в среднем во много раз превышает размеры самих частиц, соответствует только модели строения газов. В жидком и твердом агрегатных состояниях частицы вещества «упакованы» гораздо плотнее.

Ответ: $1$.
}

% Source sdamgia-3464: https://phys-ege.sdamgia.ru/problem?id=3464
\BankTask{sdamgia-3464}{
В молекулярной физике используется понятие «идеальный газ». Это понятие применимо тогда, когда можно пренебречь:

А. Средней кинетической энергией поступательного движения атомов или молекул.

Б. Энергией взаимодействия атомов или молекул, то есть действием сил притяжения или отталкивания.

В. Массой атомов или молекул.

Какое(-ие) из утверждений правильно(-ы)?

$1$) только А

$2$) только Б

$3$) только В

$4$) А и В
}{
\smallskip\noindent\textbf{Решение}\par
Идеальный газ — это математическая модель газа, в которой предполагается, что потенциальной энергией взаимодействия молекул (или атомов) можно пренебречь по сравнению с их кинетической энергией. В этой модели предполагается, что между молекулами не действуют силы притяжения или отталкивания, соударения частиц между собой и со стенками сосуда абсолютно упруги.

Таким образом, утверждение Б верно, утверждение А неверно. Поскольку кинетической энергией в модели идеального газа не пренебрегают, то не могут пренебрегать и массой частиц газа, так как кинетическая энергия этой массе пропорциональна. Следовательно, для применимости модели идеального газа необходимо выполнение только требования Б.

Ответ: $2$.
}

% Source sdamgia-3543: https://phys-ege.sdamgia.ru/problem?id=3543
\BankTask{sdamgia-3543}{
Какие положения включает модель идеального газа?

А. Объемом, занимаемым молекулами газа, можно пренебречь по сравнению с объемом сосуда.

Б. Взаимодействием молекул можно пренебречь.

В. Молекулы не сталкиваются друг с другом при движении.

$1$. А, Б и В

$2$. А и Б

$3$. А и В

$4$. Б и В
}{
\smallskip\noindent\textbf{Решение}\par
Идеальный газ — математическая модель газа, в которой предполагается, что потенциальной энергией взаимодействия молекул можно пренебречь по сравнению с их кинетической энергией. В этой модели считается, что между молекулами не действуют силы притяжения или отталкивания, то есть взаимодействием молекул пренебрегают. Модель идеального газа хорошо работает только для не очень плотных газов, поскольку в этой модели пренебрегают объемом, занимаемым молекулами газа по сравнению с объемом сосуда. А вот утверждение В к этой модели не относится. Поскольку именно за счет столкновений молекул, в результате которых они обмениваются кинетическими энергиями, газ приходит к состоянию термодинамического равновесия.

Ответ: $2$.
}

% Source sdamgia-4118: https://phys-ege.sdamgia.ru/problem?id=4118
\BankTask{sdamgia-4118}{
В учебнике по физике в одном из абзацев написано: «Молекулы считаются материальными точками, которые хаотически движутся и абсолютно упруго соударяются друг с другом и со стенками сосуда. В промежутках между столкновениями молекулы друг с другом и со стенками сосуда не взаимодействуют». Какая физическая модель описывается в этом абзаце учебника?

$1$) монокристаллическое твердое тело

$2$) поликристаллическое твердое тело

$3$) идеальная жидкость

$4$) идеальный газ
}{
\smallskip\noindent\textbf{Решение}\par
Данное описание относится к модели идеального газа.

Ответ: $4$.
}


% Source sdamgia-3297: https://phys-ege.sdamgia.ru/problem?id=3297
\BankTask{sdamgia-3297}{
При повышении температуры газа в запаянном сосуде давление газа увеличивается. Это изменение давления объясняется тем, что

$1$) увеличивается объем сосуда за счет нагревания его стенок

$2$) увеличивается энергия теплового движения молекул газа

$3$) увеличиваются размеры молекул газа при его нагревании

$4$) увеличивается энергия взаимодействия молекул газа друг с другом
}{
\smallskip\noindent\textbf{Решение}\par
Молекулы газа непрерывно ударяются о стенки сосуда, оказывая на них давление. Чем больше средняя кинетическая энергия теплового движения молекул, тем быстрее молекулы двигаются и тем больший импульс передают стенкам сосуда в единицу времени, оказывая тем самым большее давление. С другой стороны, температура является мерой энергии теплового движения молекул: $\overline{E}= \frac{3}{2} kT.$ Отсюда заключаем, что при повышении температуры газа в запаянном сосуде давление газа увеличивается вследствие увеличения энергии теплового движения молекул. 

Ответ: $2$.
}

% Source sdamgia-6372: https://phys-ege.sdamgia.ru/problem?id=6372
\BankTask{sdamgia-6372}{
Учитель записал на доске три утверждения, относящиеся к молекулам.

$1$. Размерами молекул можно пренебречь.

$2$. Молекулы при столкновениях взаимодействуют как упругие шары.

$3$. При любом расстоянии между молекулами между ними действуют силы притяжения.

Какие из этих утверждений нельзя соотнести с моделью идеального газа?

$1$) все три утверждения

$2$) только третье утверждение

$3$) первое и второе утверждения

$4$) первое и третье утверждения
}{
\smallskip\noindent\textbf{Решение}\par
В модели идеального газа предполагается, что частицы газа малы по сравнению со средним расстоянием между ними, то есть размерами частиц газа пренебрегаем. Частицы взаимодействуют только при столкновения, то есть сил притяжения между ними нет. При этом в столкновениях частицы взаимодействую как упругие шары. Таким образом, верны первое и второе утверждения. То есть утверждение под номером три нельзя соотнести с моделью идеального газа, значит, правильный ответ указан под номером $2$.

Правильный ответ указан под номером $2$.
}

% Source sdamgia-5605: https://phys-ege.sdamgia.ru/problem?id=5605
\BankTask{sdamgia-5605}{
Частицы газа находятся в среднем на таких расстояниях друг от друга, при которых силы притяжения между ними незначительны. Это объясняет

$1$) большую скорость частиц газа

$2$) значение скорости звука в газе

$3$) распространение в газе звуковых волн

$4$) способность газов к неограниченному расширению
}{
\smallskip\noindent\textbf{Решение}\par
Поскольку силы притяжения между частицами газа незначительны, то они свободно могут удаляться на любые расстояния, а значит, неограниченно расширяться.

Правильный ответ указан под номером $4$.
}

\end{document}
